% MR55--0C
% Kollaps/Synkope
% 14.0
% 2013-11-30

\label{m:kollaps}\EE{Taktik}:
\Speech{Ausführliche Untersuchung und symptomatische Behandlung 
\Ca{Verwechslung mit lebensbedrohlichen Differentialdiagnosen möglich!}}

\begin{itemize}
  \item Vitale Bedrohung einschätzen, auf \EE{Differentialdiagnosen} untersuchen, z.\,B.:
  \begin{compactitem}
    \item Insult
    \item Herzrhytmusstörungen
    \item Akutes Koronarsyndrom (evtl. schmerzlos!)
    \item Krampfanfall
    \item Schädel-Hirn-Trauma (SHT)
    \item Zuckerstoffwechselstörung (Hyper-, Hypoglykämie)
    \item Exsikkose
    \item \ldots
  \end{compactitem}
  Es sind dabei alle zur Verfügung stehenden
  Möglichkeiten auszuschöpfen, z.\,B.:
  \begin{compactitem}
    \item Traumacheck (Sturz?),
    \item Neurocheck inkl. Blutzuckermessung,
    \item Temperatur,
    \item Ausführliche Anamnese bzw. Fremdanamnese 
    \item \ldots
  \end{compactitem}
  \item \EE{Ursachenforschung} (Begleiterkrankungen (Infekt, \ldots), letzte Mahlzeit, Hitze, \ldots)
  \item Lagerung: Nach Ausschluss von Herzbeschwerden,
  Atemnot und relevanten Verletzungen: Beine hoch
  \item Transportentscheidung: Hospitalisierung zur Abklärung
  anstreben\newline
  Abteilung: Innere Medizin
\end{itemize}